\documentclass[prl,twocolumn]{revtex4}%
\usepackage{amsfonts}
\usepackage{amsmath}
\usepackage{amssymb}
\usepackage{graphicx}%
\setcounter{MaxMatrixCols}{30}
%TCIDATA{OutputFilter=latex2.dll}
%TCIDATA{Version=4.00.0.2321}
%TCIDATA{CSTFile=revtex4.cst}
%TCIDATA{Created=Thursday, January 11, 2007 19:35:53}
%TCIDATA{LastRevised=Saturday, February 10, 2007 22:06:52}
%TCIDATA{<META NAME="GraphicsSave" CONTENT="32">}
%TCIDATA{<META NAME="DocumentShell" CONTENT="Articles\SW\REVTeX 4">}
\newtheorem{theorem}{Theorem}
\newtheorem{acknowledgement}[theorem]{Acknowledgement}
\newtheorem{algorithm}[theorem]{Algorithm}
\newtheorem{axiom}[theorem]{Axiom}
\newtheorem{claim}[theorem]{Claim}
\newtheorem{conclusion}[theorem]{Conclusion}
\newtheorem{condition}[theorem]{Condition}
\newtheorem{conjecture}[theorem]{Conjecture}
\newtheorem{corollary}[theorem]{Corollary}
\newtheorem{criterion}[theorem]{Criterion}
\newtheorem{definition}[theorem]{Definition}
\newtheorem{example}[theorem]{Example}
\newtheorem{exercise}[theorem]{Exercise}
\newtheorem{lemma}[theorem]{Lemma}
\newtheorem{notation}[theorem]{Notation}
\newtheorem{problem}[theorem]{Problem}
\newtheorem{proposition}[theorem]{Proposition}
\newtheorem{remark}[theorem]{Remark}
\newtheorem{solution}[theorem]{Solution}
\newtheorem{summary}[theorem]{Summary}
\newenvironment{proof}[1][Proof]{\noindent\textbf{#1.} }{\ \rule{0.5em}{0.5em}}
\begin{document}
\preprint{HEP/123-qed}
\title[Short title for running header]{Long title that appears on the title page}
\author{First Author}
\affiliation{My Institution}
\author{Second Author}
\affiliation{My Institution}
\author{Third Author}
\affiliation{Other Institution}
\keywords{one two three}
\pacs{PACS number}

\begin{abstract}
Shell document for REV\TeX{} 4.

\end{abstract}
\volumeyear{year}
\volumenumber{number}
\issuenumber{number}
\eid{identifier}
\date[Date text]{date}
\received[Received text]{date}

\revised[Revised text]{date}

\accepted[Accepted text]{date}

\published[Published text]{date}

\startpage{101}
\endpage{102}
\maketitle
\tableofcontents


\section{}%

\[
\rho\left(  \mathbf{r}\right)  =\left\langle \left.  \Psi\right\vert
\Phi^{\dagger}\left(  \mathbf{r}\right)  \Phi\left(  \mathbf{r}\right)
\Psi\right\rangle ;\quad\mathbf{r\in}\mathbb{R}^{3}%
\]%
\begin{align*}
&  \left[  \rho\right]  _{N}=\left\{  \left.  \overset{}{\Psi}\right\vert
\Xi\left(  \Psi\right)  =\rho,\left\langle \left.  \Psi\right\vert
\Psi\right\rangle =1,\Psi\in\mathcal{H}^{N}\right. \\
&  \qquad\qquad\left.  \rho\in L_{1}\left(  \mathbb{R}^{3}\right)  \underset
{}{}\right\}  ;
\end{align*}
%

\[
\left[  \Xi\left(  \Psi\right)  \right]  \left(  \mathbf{r}\right)
=\left\langle \left.  \Psi\right\vert \Phi^{\dagger}\left(  \mathbf{r}\right)
\Phi\left(  \mathbf{r}\right)  \Psi\right\rangle ;\quad\mathbf{r\in}%
\mathbb{R}^{3}%
\]
%

\[
\left\langle \cdot|\cdot\right\rangle :L_{\infty}\left(  \mathbb{R}%
^{3}\right)  \times L_{1}\left(  \mathbb{R}^{3}\right)  \rightarrow\mathbb{R}%
\]
%

\begin{align*}
\int\rho\left(  \mathbf{r}\right)  d\mathbf{r}  &  =\left\langle \left.
\Psi\right\vert \int\Phi^{\dagger}\left(  \mathbf{r}\right)  \Phi\left(
\mathbf{r}\right)  d\mathbf{r}\Psi\right\rangle \\
&  =\left\langle \left.  \Psi\right\vert \mathcal{N}\Psi\right\rangle
=N\left\langle \left.  \Psi\right\vert \Psi\right\rangle
\end{align*}
%

\begin{align*}
&  \left[  \rho\right]  _{N}=\left\{  \left.  \overset{}{\Psi}\right\vert
\Xi\left(  \Psi\right)  =\rho,\left\langle \left.  \Psi\right\vert
\Psi\right\rangle =1,\Psi\in\mathcal{H}^{N}\right. \\
&  \qquad\qquad\left.  \rho\in L_{1}\left(  \mathbb{R}^{3}\right)  \cap
L_{3}\left(  \mathbb{R}^{3}\right)  \underset{}{}\right\}  ;
\end{align*}%
\[
\left\langle \cdot|\cdot\right\rangle :\left[  L_{\infty}\left(
\mathbb{R}^{3}\right)  \oplus L_{\frac{3}{2}}\left(  \mathbb{R}^{3}\right)
\right]  \times\left[  L_{1}\left(  \mathbb{R}^{3}\right)  \cap L_{3}\left(
\mathbb{R}^{3}\right)  \right]  \rightarrow\mathbb{R}%
\]%
\[
\rho\in L_{1}\left(  \mathbb{R}^{3}\right)  \cap L_{3}\left(  \mathbb{R}%
^{3}\right)  \equiv\mathcal{D}_{1}%
\]
%

\begin{align*}
\int\rho\left(  \mathbf{r}\right)  d\mathbf{r}  &  =\left\langle
\Psi\left\vert \int\Phi^{\dagger}\left(  \mathbf{r}\right)  \Phi\left(
\mathbf{r}\right)  d\mathbf{r}\Psi\right.  \right\rangle \\
&  =\left\langle \left.  \Psi\right\vert \mathcal{N}\Psi\right\rangle
=N\left\langle \left.  \Psi\right\vert \Psi\right\rangle
\end{align*}
%

\[
\nabla_{\eta}\mathcal{L}\left(  \Psi,\lambda,\gamma\right)  =0\equiv
\left\langle \eta\left\vert \gamma\right.  \right\rangle -\left\langle \left.
\Psi\right\vert \Omega\left(  \eta\right)  \Psi\right\rangle =0
\]%
\[
\forall\eta\in L_{\infty}\left(  \mathbb{R}^{3}\right)  \oplus L_{\frac{3}{2}%
}\left(  \mathbb{R}^{3}\right)
\]


$\Psi\in$ Domain of $H_{U}$ $\Leftrightarrow$ $\Psi\in$ Domain of
$\Omega\left(  \lambda\right)  $
\begin{align*}
\nabla_{\eta\Phi}^{2}\mathcal{L}\left(  \Psi\left(  \rho\right)
,\lambda\left(  \rho\right)  ,\rho\right)   &  =0\,\forall\eta\in L_{\infty
}\left(  \mathbb{R}^{3}\right)  \oplus L_{\frac{3}{2}}\left(  \mathbb{R}%
^{3}\right) \\
&  \Updownarrow\\
\left\langle \left.  \Psi\right\vert \Omega\left(  \eta\right)  \Phi
\right\rangle  &  =0\,\forall\eta\in L_{\infty}\left(  \mathbb{R}^{3}\right)
\oplus L_{\frac{3}{2}}\left(  \mathbb{R}^{3}\right)
\end{align*}
%

\[
\left\{  \left.  \left\vert \Psi_{\kappa}\left(  \rho\right)  \right\rangle
\right\vert \quad1\leq\kappa\leq\tau\right\}
\]%
\begin{align*}
\mathcal{F}_{HK}\left(  \rho\right)   &  =\left\langle \Psi_{\kappa}\left(
\rho\right)  \left\vert H_{U}\right.  \Psi_{\kappa}\left(  \rho\right)
\right\rangle \\
\left\vert \Psi_{\mathbf{c}}\left(  \rho\right)  \right\rangle  &
\in\text{linear span}\left\{  \left.  \left\vert \Psi_{\kappa}\left(
\rho\right)  \right\rangle \right\vert \quad1\leq\kappa\leq\tau\right\}
\end{align*}%
\[
\left\{  \left.  \left\vert \Psi_{\kappa}\left(  \rho\right)  \right\rangle
\right\vert \quad1\leq\kappa\leq\tau\right\}
\]%
\[
\rho\in L_{1}\left(  \mathbb{R}^{3}\right)  \cap L_{3}\left(  \mathbb{R}%
^{3}\right)  \equiv\mathcal{D}_{1}%
\]%
\[
\left\langle \cdot|\cdot\right\rangle :\mathcal{D}_{1}^{\ast}\times
\mathcal{D}_{1}\rightarrow\mathbb{R}%
\]%
\[
\mathcal{D}_{1}^{\ast}=\left[  L_{\infty}\left(  \mathbb{R}^{3}\right)  \oplus
L_{\frac{3}{2}}\left(  \mathbb{R}^{3}\right)  \right]
\]%
\[
\mathcal{D}\left(  H_{U}\right)
\]%
\[
\mathcal{D}_{1}^{\ast}=\left[  L_{\infty}\left(  \mathbb{R}^{3}\right)  \oplus
L_{\frac{3}{2}}\left(  \mathbb{R}^{3}\right)  \right]
\]%
\[
\mathcal{L}\left(  \Psi\left(  \rho\right)  ,\lambda\left(  \rho\right)
\right)  =\inf_{\Psi\in\mathcal{D}\left(  H_{U}\right)  }\sup_{\lambda
}\mathcal{L}\left(  \Psi,\lambda,\rho\right)
\]%
\[
\left\{  \left.  \Omega\left(  \lambda\right)  \right\vert \lambda
\in\mathcal{D}_{1}^{\ast}\right\}
\]%
\[
H_{U}%
\]%
\[
=\left\langle \left.  \Psi\left(  \rho\right)  \right\vert \Omega\left(
\lambda\right)  \Psi\left(  \rho\right)  \right\rangle
\]
1%
\[
H\left(  \lambda\right)  =\sum_{1\leq j\leq\kappa}E_{\kappa}\left\vert
\Lambda_{\kappa}\right\rangle \left\langle \Lambda_{\kappa}\right\vert
+\int_{\sigma_{c}\left(  H\left(  \lambda\right)  \right)  }E\left\vert
\Lambda\left(  E\right)  \right\rangle \left\langle \Lambda\left(  E\right)
\right\vert dE
\]
2%
\[
\mathcal{P}_{d}\left(  H\left(  \lambda\right)  \right)  =\sum_{1\leq
j\leq\kappa}\left\vert \Lambda_{\kappa}\right\rangle \left\langle
\Lambda_{\kappa}\right\vert
\]
3%
\[
\mathcal{P}_{c}\left(  H\left(  \lambda\right)  \right)  =\int_{\sigma
_{c}\left(  H\left(  \lambda\right)  \right)  }\left\vert \Lambda\left(
E\right)  \right\rangle \left\langle \Lambda\left(  E\right)  \right\vert dE
\]
4%
\[
\mathcal{H}^{N}=\mathcal{H}_{d}^{N}\oplus\mathcal{H}_{c}^{N}%
\]
5%
\[
\mathcal{H}_{d}^{N}=\mathcal{P}_{d}\left(  H\left(  \lambda\right)  \right)
\mathcal{H}^{N};\qquad\mathcal{H}_{c}^{N}=\mathcal{P}_{c}\left(  H\left(
\lambda\right)  \right)  \mathcal{H}^{N}%
\]
6%
\[
H\left(  \lambda\right)
\]
7%
\[
H_{U}%
\]%
\begin{align*}
&  \left.  \qquad\right.  \left.  \qquad\right.  \mathcal{L}\left(  \Psi
+\Phi,\lambda+\eta,\rho\right) \\
&  \left.  \qquad\right.  \left.  \qquad\right.  \left.  \qquad\right.
\parallel\\
&  \left\langle \left(  \Psi+\Phi\right)  \left\vert H\left(  \lambda\right)
\right.  \left(  \Psi+\Phi\right)  \right\rangle +\left\langle \left(
\lambda+\eta\right)  \left\vert \rho\right.  \right\rangle \\
&  \left.  \qquad\right.  \left.  \qquad\right.  \left.  \qquad\right.
\parallel\\
&  \quad\left\langle \Psi\left\vert H\left(  \lambda\right)  \right.
\Psi\right\rangle +\left\langle \lambda\left\vert \rho\right.  \right\rangle
+\left\langle \Phi\left\vert H\left(  \lambda\right)  \right.  \Psi
\right\rangle \\
&  +\left\langle \Psi\left\vert H\left(  \lambda\right)  \right.
\Phi\right\rangle +\left\langle \Phi\left\vert H\left(  \lambda\right)
\right.  \Phi\right\rangle +\left\langle \eta\left\vert \rho\right.
\right\rangle
\end{align*}%
\[
\left(  \Psi,\lambda\right)
\]%
\[
\frac{\sigma_{\inf}\left(  \rho\right)  }{N}\int\Phi^{\dagger}\left(
\mathbf{r}\right)  \Phi\left(  \mathbf{r}\right)  d\mathbf{r}%
\]%
\[
\sigma_{\inf}\left(  \rho\right)  =\inf\limits_{\Psi}\left\{  \left\langle
\Psi\left\vert H\left(  \lambda\right)  \right.  \Psi\right\rangle \right\}
\]%
\[
\Omega\left(  \lambda\right)  \rightarrow\Omega\left(  \tilde{\lambda}\right)
=\int\tilde{\lambda}\left(  \mathbf{r}\right)  \Phi^{\dagger}\left(
\mathbf{r}\right)  \Phi\left(  \mathbf{r}\right)  d\mathbf{r=}\int\left(
\lambda\left(  \mathbf{r}\right)  -\sigma_{\inf}\left(  \rho\right)  \right)
\Phi^{\dagger}\left(  \mathbf{r}\right)  \Phi\left(  \mathbf{r}\right)
d\mathbf{r}%
\]
Normalization constraint%
\[
\mathfrak{N}\left(  \Psi\right)  =\int\left\langle \left.  \Psi\right\vert
\Phi^{\dagger}\left(  \mathbf{r}\right)  \Phi\left(  \mathbf{r}\right)
\Psi\right\rangle d\mathbf{r}=\int\rho\left(  \mathbf{r}\right)
d\mathbf{r}=N
\]%
\[
\mathfrak{N}\left(  \Psi\right)  =\int\left\langle \left.  \Psi\right\vert
\Phi^{\dagger}\left(  \mathbf{r}\right)  \Phi\left(  \mathbf{r}\right)
\Psi\right\rangle d\mathbf{r}%
\]%
\[
\mathfrak{N}\left(  \Psi\right)  =N
\]
%

\[
\mathcal{L}\left(  \Psi,\lambda,\rho\right)  =\left\langle \left.
\Psi\right\vert H_{U}\Psi\right\rangle -\left\langle \lambda\left\vert
\left\{  \overset{}{\Xi\left(  \Psi\right)  -\rho}\right\}  \right.
\right\rangle -\beta\left(  \mathfrak{N}\left(  \Psi\right)  -N\right)
\]
leading to
\[
\nabla\mathcal{L}\left(  \bar{\Psi},\Psi,\lambda,\rho\right)  =0=\nabla
E_{U}\left(  \bar{\Psi},\Psi\right)  -\int\lambda\left(  \mathbf{r}\right)
\nabla\Xi\left(  \mathbf{r}\right)  \left(  \bar{\Psi},\Psi\right)
d\mathbf{r-}\beta\nabla\mathfrak{N}\left(  \Psi\right)
\]%
\[
\mathcal{L}\left(  \Psi,\lambda,\rho\right)  =\left\langle \left.
\Psi\right\vert H\left(  \lambda+\beta\right)  \Psi\right\rangle +\left\langle
\left.  \lambda\right\vert \rho\right\rangle +\beta N
\]%
\[
\int\left\langle \left.  \Psi\right\vert \Phi^{\dagger}\left(  \mathbf{r}%
\right)  \Phi\left(  \mathbf{r}\right)  \Psi\right\rangle d\mathbf{r}=N
\]%
\[
\sum_{1\leq j\leq m}\left\langle \left.  \Psi\right\vert \Phi_{j}^{\dagger
}\Phi_{j}\Psi\right\rangle =N
\]%
\begin{equation}
\Xi\left(  \Psi\right)  =\left(  \Xi_{1}\left(  \Psi\right)  ,\cdots,\Xi
_{r}\left(  \Psi\right)  \right)
\end{equation}%
\begin{equation}
\Xi_{j}\left(  \Psi\right)  =\left\langle \left.  \Psi\right\vert \Phi
_{j}^{\dagger}\Phi_{j}\Psi\right\rangle
\end{equation}%
\begin{equation}
\left\langle \left.  \Psi\right\vert H_{U}\Psi\right\rangle =\sum
_{\substack{1\leq j_{1}<\cdots<j_{n}\leq m\\1\leq k_{1}<\cdots<k_{n}\leq
m}}\bar{c}_{j_{1}\cdots j_{n}}c_{k_{1}\cdots k_{n}}\left[  H_{U}\right]
_{j_{1}\cdots j_{n}k_{1}\cdots k_{n}}%
\end{equation}%
\begin{equation}
\nabla_{\bar{\Psi}\Psi}\mathcal{L}\left(  \bar{\Psi},\Psi,\lambda,\rho\right)
\equiv\nabla_{\bar{\Psi}\Psi}\mathcal{L}\left(  \mathbf{\bar{C}}%
,\mathbf{C},\lambda,\rho\right)  \equiv\left(  \nabla_{\mathbf{\bar{C}}%
}\mathcal{L}\left(  \mathbf{\bar{C}},\mathbf{C},\lambda,\rho\right)
,\nabla_{\mathbf{C}}\mathcal{L}\left(  \mathbf{\bar{C}},\mathbf{C}%
,\lambda,\rho\right)  \right)
\end{equation}%
\begin{equation}
\nabla_{\mathbf{C}}\mathcal{L}\left(  \mathbf{\bar{C}},\mathbf{C},\lambda
,\rho\right)  =\left(  \frac{\partial\mathcal{L}\left(  \mathbf{\bar{C}%
},\mathbf{C},\lambda,\rho\right)  }{\partial c_{1\cdots n}},\cdots
,\frac{\partial\mathcal{L}\left(  \mathbf{\bar{C}},\mathbf{C},\lambda
,\rho\right)  }{\partial c_{m-n\cdots m}}\right)
\end{equation}%
\begin{equation}
\nabla_{\mathbf{\Psi}}\Xi_{1}\left(  \bar{\Psi},\Psi\right)  \equiv\nabla
\Xi_{1}\left(  \mathbf{\bar{C}},\mathbf{C}\right)  =\left(  \frac{\partial
\Xi_{1}\left(  \mathbf{\bar{C}},\mathbf{C}\right)  }{\partial c_{1\cdots n}%
},\cdots,\frac{\partial\Xi_{1}\left(  \mathbf{\bar{C}},\mathbf{C}\right)
}{\partial c_{m-n\cdots m}}\right)
\end{equation}%
\begin{equation}
\nabla_{\mathbf{\Psi}}\mathfrak{N}\left(  \bar{\Psi},\Psi\right)  \equiv
\nabla\mathfrak{N}\left(  \mathbf{\bar{C}},\mathbf{C}\right)  =\left(
\frac{\partial\mathfrak{N}\left(  \mathbf{\bar{C}},\mathbf{C}\right)
}{\partial c_{1\cdots n}},\cdots,\frac{\partial\mathfrak{N}\left(
\mathbf{\bar{C}},\mathbf{C}\right)  }{\partial c_{m-n\cdots m}}\right)
\end{equation}%
\begin{equation}
\nabla_{\bar{\Psi}}\mathcal{L}\left(  \bar{\Psi},\Psi,\mathbf{\lambda}%
,\rho\right)  =0=\left\langle \left.  {}\right\vert H\left(  \mathbf{\lambda
}+\beta\right)  \Psi\right\rangle
\end{equation}%
\begin{equation}
\nabla_{\left(  \mathbf{\lambda},\beta\right)  }\mathcal{L}\left(  \bar{\Psi
},\Psi,\mathbf{\lambda},\rho\right)  =0=\left\langle \left.  {}\right\vert
\rho\right\rangle +N
\end{equation}
Find $\left(  \mathbf{\lambda},\beta\right)  $ such that one can find
solutions $\left(  \bar{\Psi},\Psi\right)  .$
\begin{equation}
\left\{  H_{U}-\Omega\left(  \lambda+\beta\right)  \right\}  \left\vert
\Psi\right\rangle =0
\end{equation}
%

\[
\mathcal{L}\left(  \Psi,\lambda,\rho\right)  =\left\langle \left.
\Psi\right\vert H_{U}\Psi\right\rangle -\left\langle \lambda\left\vert
\left\{  \overset{}{\Xi\left(  \Psi\right)  -\rho}\right\}  \right.
\right\rangle -\beta\left\{  \mathfrak{N}\left(  \Psi\right)  -N\right\}
\]%
\[
\rho
\]%
\begin{align*}
&  \left\langle \lambda\left\vert \overset{}{\left\{  \Xi\left(  \Psi\right)
-\rho\right\}  }\right.  \right\rangle +\beta\left\{  \mathfrak{N}\left(
\Psi\right)  -N\right\} \\
&  \qquad\qquad\qquad\qquad\left.  \shortparallel\right. \\
&  \left\langle \left.  \Psi\right\vert \Omega\left(  \lambda+\beta\right)
\Psi\right\rangle -\left\langle \lambda\left\vert \rho\right.  \right\rangle
-\beta N
\end{align*}
%

\[
\Omega\left(  \lambda\right)  =\int\left\{  \lambda\left(  \mathbf{r}\right)
+\beta\right\}  \mathbf{\Phi}\left(  \mathbf{r}\right)  ^{\dagger}%
\mathbf{\Phi}\left(  \mathbf{r}\right)  d\mathbf{r}%
\]
%

\begin{align*}
&  \qquad\qquad\qquad\qquad\mathcal{L}\left(  \Psi,\lambda,\rho\right) \\
&  \qquad\qquad\qquad\qquad\qquad\left.  \shortparallel\right. \\
&  \left\langle \left.  \Psi\right\vert \left\{  K+V_{2}-\Omega\left(
\lambda+\beta\right)  \right\}  \Psi\right\rangle +\left\langle \lambda
\left\vert \rho\right.  \right\rangle +\beta N
\end{align*}
%

\[
\left\{  \left.  \Omega\left(  \lambda+\beta\right)  \right\vert \lambda
\in\mathcal{D}_{1}^{\ast};\beta\in\mathbb{R}\right\}
\]
%

\[
H\left(  \lambda\right)  =K+V_{2}-\Omega\left(  \lambda+\beta\right)
\]
%

\[
H\left(  \lambda\right)  =K+V_{2}-\Omega\left(  \lambda+\beta\right)
\]%
\begin{align*}
&  \qquad\qquad\qquad\qquad\mathcal{L}\left(  \Psi,\lambda,\beta,\rho\right)
\\
&  \qquad\qquad\qquad\qquad\qquad\left.  \shortparallel\right. \\
&  \left\langle \left.  \Psi\right\vert H_{U}\Psi\right\rangle -\left\langle
\lambda\left\vert \left\{  \overset{}{\Xi\left(  \Psi\right)  -\rho}\right\}
\right.  \right\rangle -\beta\left\{  \mathfrak{N}\left(  \Psi\right)
-N\right\}
\end{align*}
%

\begin{align*}
&  \left[  \rho\right]  _{N}=\left\{  \left.  \overset{}{\Psi}\right\vert
\Xi\left(  \Psi\right)  =\rho,\mathfrak{N}\left(  \Psi\right)  =N,\Psi
\in\mathcal{H}^{N}\right. \\
&  \qquad\qquad\left.  \rho\in L_{1}\left(  \mathbb{R}^{3}\right)  \cap
L_{3}\left(  \mathbb{R}^{3}\right)  \underset{}{}\right\}  ;
\end{align*}%
\[
\mathcal{L}\left(  \Psi\left(  \rho\right)  ,\lambda\left(  \rho\right)
\right)  =\inf_{\Psi\in\mathcal{D}\left(  H_{U}\right)  }\sup_{\lambda,\beta
}\mathcal{L}\left(  \Psi,\lambda,\rho\right)
\]
%

\[
\mathcal{L}\left(  \Psi\left(  \rho\right)  ,\lambda\left(  \rho\right)
\right)  =\inf_{\Psi\in\mathcal{D}\left(  H_{U}\right)  }\sup_{\lambda,\beta
}\mathcal{L}\left(  \Psi,\lambda,\rho\right)
\]


$\beta\in\mathbb{R}$ is a free parameter, while $\lambda\in\mathcal{D}%
_{1}^{\ast}$ is not

i.e. one must find specific optimal values for $\lambda$ in

order to find a constrained inf w.r.t $\Psi$, while any

value of $\beta$ suffices. Thus one can always shift the

optimal value of $\lambda$ by a constant and still obtain the

the same inf over $\Psi$%

\[
H\left(  \lambda\right)  =K+V_{2}-\Omega\left(  \lambda\right)
\]
%

\[
\mathcal{L}\left(  \Psi\left(  \rho\right)  ,\lambda\left(  \rho\right)
\right)  =\inf_{\Psi\in\mathcal{D}\left(  H_{U}\right)  }\sup_{\lambda
}\mathcal{L}\left(  \Psi,\lambda,\rho\right)
\]%
\begin{align*}
\left.
\begin{array}
[c]{c}%
\lambda\left(  \rho\right) \\
\Psi\left(  \rho\right)
\end{array}
\right\}  \text{ exist }  &  \Longleftrightarrow\text{V-representability}\\
\lambda\left(  \rho\right)  \text{ is unique }  &  \Longleftrightarrow
\text{Differentialbility}\\
\Psi\left(  \rho\right)  \text{ is unique }  &  \Longleftrightarrow
\text{Degeneracy}%
\end{align*}


\quad If for a given $\rho$ there exists a Hamiltonian $H\left(
\tilde{\lambda}\left(  \rho\right)  \right)  $

that has at least one eigen value that is lower than the

infinum of its continuous spectra then there exists a

state $\left\vert \Psi\left(  \rho\right)  \right\rangle $ that satisfies the
First Order Necessary

$\qquad$Conditions for a local potential $\Omega\left(  \lambda\left(
\rho\right)  \right)  ,$

\qquad$\qquad\qquad$where $\lambda\left(  \rho\right)  -\tilde{\lambda}\left(
\rho\right)  =const.$

$\left\vert \Psi\left(  \rho\right)  \right\rangle $
\begin{align*}
&  \text{Let }\Phi\left(  \mathbf{r}\right)  ^{\dagger}\Phi\left(
\mathbf{r}\right)  \left\vert \Psi\right\rangle \overset{a.e}{\neq
}0;\,\mathbf{r\in}\operatorname*{supp}\nolimits_{1}\left\{  \Psi\right\}
\mathbf{\subseteq}\mathbb{R}^{3}\\
&  \qquad\qquad\text{then }\Omega\left(  \lambda\right)  \left\vert
\Psi\right\rangle =0\Leftrightarrow\lambda\overset{a.e}{=}0.
\end{align*}


%

\begin{align*}
&  \text{By a.e. we mean almost every where in}\operatorname*{supp}%
\nolimits_{1}\left\{  \Psi\right\}  \text{, where}\\
&  \operatorname*{supp}\nolimits_{1}\left\{  \Psi\right\}
=\operatorname*{closure}\left\{  \mathbf{r}\left\vert \Psi\left(
\mathbf{r},\mathbf{r}_{2},\cdots,\mathbf{r}_{N}\right)  \neq0;\right.
\right.  \\
&  \qquad\qquad\qquad\qquad\qquad\left.  \mathbf{r\neq r}_{j}\neq
\mathbf{r}_{k}\in\mathbb{R}^{3},2\leq j\neq k\leq N\right\}
\end{align*}
1

$\mathcal{B}_{1}\left(  \mathcal{H}^{N}\right)  $ is the space of Trace Class
operators
\[
\mathcal{B}_{1}\left(  \mathcal{H}^{N}\right)  =\left\{  X\left\vert
\operatorname{Tr}\left\{  \left(  X^{\dagger}X\right)  ^{\frac{1}{2}}\right\}
<\infty\right.  \right\}
\]
2

$\mathcal{B}_{2}\left(  \mathcal{H}^{N}\right)  $ is the space of Hilbert
Schimdt operators

3%
\[
\mathcal{B}_{2}\left(  \mathcal{H}^{N}\right)  =\left\{  X\left\vert \left\{
\operatorname{Tr}\left\{  X^{\dagger}X\right\}  \right\}  ^{\frac{1}{2}%
}<\infty\right.  \right\}
\]
4%
\[
\mathcal{B}_{2}\left(  \mathcal{H}^{N}\right)  \text{is a Hilbert Space with
inner product }\left(  X\left\vert Y\right.  \right)  =\operatorname{Tr}%
\left\{  X^{\dagger}Y\right\}
\]
5%
\begin{align*}
\mathcal{B}_{\kappa}\left(  \mathcal{H}^{N}\right)    & =\mathcal{B}_{\kappa
sa}\left(  \mathcal{H}^{N}\right)  \oplus\mathcal{B}_{\kappa asa}\left(
\mathcal{H}^{N}\right)  \\
\mathcal{B}_{\kappa sa}\left(  \mathcal{H}^{N}\right)    & =\left\{  \left.
X=X^{\dagger}\right\vert X\in\mathcal{B}_{\kappa}\left(  \mathcal{H}%
^{N}\right)  \right\}  \\
\mathcal{B}_{\kappa asa}\left(  \mathcal{H}^{N}\right)    & =\left\{  \left.
X=-X^{\dagger}\right\vert X\in\mathcal{B}_{\kappa}\left(  \mathcal{H}%
^{N}\right)  \right\}  \\
\kappa & =1,2
\end{align*}


6%
\begin{align*}
& \quad\quad\quad\quad\text{The set of }N\text{-electron Density Operators
}D^{N}\\
& S_{N}\subset\mathcal{B}_{1sa+}\left(  \mathcal{H}^{N}\right)  \equiv
\text{The positive elemements of }\mathcal{B}_{1sa}\left(  \mathcal{H}%
^{N}\right)
\end{align*}
7%
\[
D^{N}=A^{2};A\in\mathcal{B}_{2sa}\left(  \mathcal{H}^{N}\right)
\]
8%
\[
\mathcal{B}_{2sa}\left(  \mathcal{H}^{N}\right)  \text{ is a real Hilbert
Space with symmetric inner product}%
\]
9%
\begin{align*}
& \quad\text{Define Super Operators }\\
& \hat{T}:\mathcal{B}_{2sa}\left(  \mathcal{H}^{N}\right)  \rightarrow
\mathcal{B}_{2sa}\left(  \mathcal{H}^{N}\right)  \text{ }\\
& \widehat{T}\left\vert A\right)  =\left\vert \frac{\left[  T,A\right]  _{+}%
}{2}\right)  \\
& \left[  T,Y\right]  _{+}=TA+AT
\end{align*}
10%
\begin{align*}
\Xi\left(  A\right)    & =\left(  A\left\vert \widehat{\Phi^{\dagger}\Phi
}A\right.  \right)  =\rho\\
\mathfrak{N}\left(  A\right)    & =\int\left\langle \left.  A\right\vert
\widehat{\Phi^{\dagger}\left(  \mathbf{r}\right)  \Phi\left(  \mathbf{r}%
\right)  }A\right\rangle d\mathbf{r=}N
\end{align*}
12%
\[
\mathcal{F}_{HK}\left(  \rho\right)  =\inf_{A\in\mathcal{B}_{2sa}\left(
\mathcal{H}^{N}\right)  }\operatorname{Tr}\left\{  A\overset{}{\left(
K+V_{2}\right)  A}\right\}  =\left(  A\left(  \rho\right)  \left\vert
\widehat{\left(  K+V_{2}\right)  }A\left(  \rho\right)  \right.  \right)
\]
13%
\[
\mathcal{L}\left(  A,\lambda,\gamma,\beta\right)  =\left(  A\left\vert
\widehat{\left(  K+V_{2}-\Omega\left(  \lambda+\beta\right)  \right)
}A\right.  \right)  +\left\langle \lambda\left\vert \gamma\right.
\right\rangle +\beta N
\]

begin{align*}
&  \qquad\qquad\qquad\qquad\mathcal{L}\left(  \Psi,\lambda,\beta,\rho\right)
\\
&  \qquad\qquad\qquad\qquad\qquad\left.  \shortparallel\right.  \\
&  \left\langle \left.  \Psi\right\vert H_{U}\Psi\right\rangle -\left\langle
\lambda\left\vert \left\{  \overset{}{\Xi\left(  \Psi\right)  -\rho}\right\}
\right.  \right\rangle -\beta\left\{  \mathfrak{N}\left(  \Psi\right)
-N\right\}
\end{align*}

\end{document} 